\uuid{3uRq}
\titre{Injection/surjection}
\auteur{bodin ; hanani ; mzari}
\organisation{exo7}
\licence{CC BY-NC-SA 4.0}
\datecreate{2020-01-20}
\niveau{L1}
\module{Algèbre}
\chapitre{Logique, ensemble, raisonnement}
\sousChapitre{Ensemble}
\difficulte{3}
\theme{Ensembles, applications, injection, surjection}
\exo7qcmid{lille-1.ensembles.030}

\contenu{
\begin{question}
\qclassification{100.02,101.01,102.01,102.02}
\qtags{motcle=injection/surjection}
Soit $f:[-1,1]\to [-1,1]$ l'application définie par
$$\forall x\in [-1,1],\; f(x)=\frac{2x}{1+x^2}.$$
Quelle sont les bonnes réponses ?
\begin{answers}
    \bad{$f$ est injective mais non surjective.}
    \bad{$f$ est surjective mais non injective.}
    \bad{$f$ n'est ni injective ni surjective.}
    \good{$f$ est bijective et $\displaystyle f^{-1}(x)=\frac{x}{1+\sqrt{1-x^2}}$.}
\end{answers}
\begin{explanations}
Soit $y\in [-1,1]$. On a $\displaystyle f(x)=y\Leftrightarrow yx^2-2x+y=0$. On résout dans $[-1,1]$ cette équation, d'inconnue $x$. Si $y=0$, on aura $x=0$ et si $y\neq 0$, on calcule $\Delta =4(1-y^2)\geq 0$ et donc
$$x=\frac{1-\sqrt{1-y^2}}{y}=\frac{y}{1+\sqrt{1-y^2}}\in [-1,1]\mbox{ car }\frac{1+\sqrt{1-y^2}}{y}\notin [-1,1].$$
Ainsi tout $y\in [-1,1]$ admet un unique antécédent $\displaystyle x=\frac{y}{1+\sqrt{1-y^2}}\in [-1,1]$. Donc $f$ est bijective et $\displaystyle f^{-1}(y)=\frac{y}{1+\sqrt{1-y^2}}$.
\end{explanations}
\end{question}
}
